\begin{figure}[H]
\centering
\begin{minipage}[t]{0.48\textwidth}
  \vspace{0pt}\centering
  \includegraphics[width=\linewidth]{figures/ch02/panels/fig2_2a_latent_cov.pdf}
\end{minipage}\hfill
\begin{minipage}[t]{0.48\textwidth}
  \vspace{0pt}\centering
  \includegraphics[width=\linewidth]{figures/ch02/panels/fig2_2b_noise_cov.pdf}
\end{minipage}

\vspace{0.8em}

\begin{minipage}[t]{0.48\textwidth}
  \vspace{0pt}\centering
  \includegraphics[width=\linewidth]{figures/ch02/panels/fig2_2c_total_cov.pdf}
\end{minipage}\hfill
\begin{minipage}[t]{0.48\textwidth}
  \vspace{0pt}
  \caption{\textbf{PPCA 中的协方差分解。}
  在图示的 $k=1,d=2$ 情形中，$\mathbf{W}\mathbf{W}^{\top}$ 为秩一矩阵，只沿 $\mathrm{span}(\mathbf{W})$ 贡献方差；$\alpha^2\mathbf{I}$ 则是各向同性且满秩的。二者之和给出满秩的边缘协方差，同时描述沿潜在子空间及其周围的变化。实线与虚线分别表示一个和两个标准差的范围。}
  \label{fig:ppca-covariance-decomposition}
\end{minipage}
\end{figure}
